\documentclass[12pt,a4paper]{ctexart}
\usepackage[margin=2.2cm]{geometry}
\usepackage{graphicx}
\usepackage{booktabs}
\usepackage{amsmath}
\usepackage{float}
\usepackage{hyperref}
\usepackage{xcolor}
\usepackage{listings}

\lstset{
	language=C,
	basicstyle=\ttfamily\small,
	keywordstyle=\color{blue!70!black},
	commentstyle=\color{green!40!black},
	stringstyle=\color{orange!60!black},
	showstringspaces=false,
	breaklines=true,
	frame=single,
	tabsize=4
}

% ===== 模板配置 - 请根据实验要求修改 =====
\newcommand{\reporttitle}{实验四：SLR(1)分析表生成}
\newcommand{\subtask}{SLR(1) ACTION/GOTO 分析表构造}
\newcommand{\authname}{姓名}
\newcommand{\authid}{学号}
% ==========================================

\title{\reporttitle}
\author{\authname：\underline{李祥熙} \quad \authid：\underline{2234412799}}
\date{\today}

\begin{document}
\maketitle
\tableofcontents
\newpage

\section{实验目的}
本实验的目标如下：
\begin{itemize}
	\item 目标一：\textbf{构建给定文法的 LR(0) 项目集规范族与状态转移图}
	\item 目标二：\textbf{计算文法符号的 FIRST/FOLLOW 集，为 SLR(1) 分析表提供依据}
	\item 目标三：\textbf{生成 SLR(1) 的 ACTION/GOTO 表并检测冲突}
\end{itemize}

\section{实验过程}
\subsection{实验环境}
\begin{itemize}
	\item 操作系统：\textbf{Linux}
	\item 编程语言：\textbf{C}
	\item 编译器：\textbf{GCC}
	\item 其他工具：\textbf{终端与文本编辑器}
\end{itemize}

\subsection{输入文法与预处理}
本实验的输入为文本形式的文法规则（支持 ``\texttt{|}'' 表示同一产生式的多个候选式）。
本次运行使用表达式文法（文件：\texttt{grammar\_expr.txt}）：
\begin{verbatim}
E -> E + T | T
T -> T * F | F
F -> ( E ) | id
\end{verbatim}
预处理步骤为：
\begin{itemize}
	\item 读取所有规则，识别非终结符与终结符集合。
	\item 自动增广文法，添加新的开始符号 $E'$。
	\item 将 ``\texttt{|}'' 分支拆分为多条产生式，便于后续计算。
\end{itemize}

\subsection{实现步骤}
\begin{enumerate}
	\item \textbf{文法解析与增广}：读取输入、建立符号表并生成增广产生式。
	\item \textbf{LR(0) 项目集}：计算闭包与 GOTO，构建规范族及状态转移图。
	\item \textbf{FIRST/FOLLOW}：迭代求解 FIRST 与 FOLLOW 集（处理空产生式）。
	\item \textbf{SLR(1) 表构造}：依据项目集与 FOLLOW 集填充 ACTION/GOTO 表，并检测冲突。
\end{enumerate}

\section{实验内容}
\subsection{任务一：LR(0) 项目集规范族构造}
实现要点：
\begin{itemize}
	\item 项目用 ``(产生式编号, 圆点位置)'' 表示，便于比较与排序。
	\item 通过闭包函数将 ``\texttt{A -> \alpha . B \beta}'' 推导出的项目递归加入。
	\item GOTO 基于某个符号移动圆点，再对结果求闭包。
\end{itemize}

\textbf{代码片段（闭包与 GOTO）}：
\begin{lstlisting}[language=C]
static void
closure (ItemSet *set, const Production *prods, const int *nonterm_first,
				 const int *nonterm_count, const int *symbol_is_nonterm)
{
	int changed = 1;
	while (changed)
		{
			changed = 0;
			for (int i = 0; i < set->count; ++i)
				{
					Item it = set->items[i];
					const Production *p = &prods[it.prod];
					if (it.dot >= p->rhs_len)
						{
							continue;
						}
					int sym = p->rhs[it.dot];
					if (!symbol_is_nonterm[sym])
						{
							continue;
						}
					int nt_index = sym;
					int start = nonterm_first[nt_index];
					int cnt = nonterm_count[nt_index];
					for (int k = 0; k < cnt; ++k)
						{
							Item new_item = { start + k, 0 };
							if (!itemset_contains (set, &new_item))
								{
									itemset_add (set, &new_item);
									changed = 1;
								}
						}
				}
		}
}
\end{lstlisting}
\textbf{说明}：该片段展示了闭包算法的核心逻辑：当圆点后为非终结符时，将其所有产生式加入项目集并反复迭代，直到集合稳定。

\subsection{任务二：SLR(1) ACTION/GOTO 表生成}
实现要点：
\begin{itemize}
	\item 对每个状态的项目：
		\begin{itemize}
			\item 若圆点前进符号为终结符，则填写 ACTION 的移进（shift）。
			\item 若为归约项目，则对该产生式左部的 FOLLOW 集填写归约（reduce）。
			\item 若为增广产生式的归约项目，则在 ``$'' 列填接受（accept）。
		\end{itemize}
	\item 对每个非终结符的状态转移填写 GOTO。
	\item 若同一格出现不同动作，判为 SLR(1) 冲突。
\end{itemize}

\section{实验结果}
\subsection{运行与输出}
\begin{figure}[H]
\centering
\includegraphics[width=0.95\textwidth]{images/compile_and_run.png}
\caption{编译与运行示例}
\label{fig:compile}
\end{figure}
\textbf{说明}：图 \ref{fig:compile} 展示了在实验目录下使用 GCC 编译，并执行 ``\texttt{./slr1 grammar\_expr.txt E}'' 的过程，便于验证程序能够顺利完成分析表生成。

\subsection{输出片段与分析}
程序输出包含：增广文法、LR(0) 项目集、FOLLOW 集、ACTION 表与 GOTO 表。以下是程序输出：
\begin{verbatim}
========================================
Augmented Grammar
========================================
0: E' -> E
1: E -> E + T
2: E -> T
3: T -> T * F
4: T -> F
5: F -> ( E )
6: F -> id

========================================
LR(0) Canonical Collection + Closure
========================================
State 0:
  E' -> . E
  E -> . E + T
  E -> . T
  T -> . T * F
  T -> . F
  F -> . ( E )
  F -> . id
Kernel: E' -> . E

State 1:
  E' -> E .
  E -> E . + T
Kernel: E' -> E ., E -> E . + T

State 2:
  E -> T .
  T -> T . * F
Kernel: E -> T ., T -> T . * F

State 3:
  T -> F .
Kernel: T -> F .

State 4:
  E -> . E + T
  E -> . T
  T -> . T * F
  T -> . F
  F -> . ( E )
  F -> ( . E )
  F -> . id
Kernel: F -> ( . E )

State 5:
  F -> id .
Kernel: F -> id .

State 6:
  E -> E + . T
  T -> . T * F
  T -> . F
  F -> . ( E )
  F -> . id
Kernel: E -> E + . T

State 7:
  T -> T * . F
  F -> . ( E )
  F -> . id
Kernel: T -> T * . F

State 8:
  E -> E . + T
  F -> ( E . )
Kernel: E -> E . + T, F -> ( E . )

State 9:
  E -> E + T .
  T -> T . * F
Kernel: E -> E + T ., T -> T . * F

State 10:
  T -> T * F .
Kernel: T -> T * F .

State 11:
  F -> ( E ) .
Kernel: F -> ( E ) .

========================================
State Transition Graph
========================================
0 --E--> 1
0 --T--> 2
0 --F--> 3
0 --(--> 4
0 --id--> 5
1 --+--> 6
2 --*--> 7
4 --E--> 8
4 --T--> 2
4 --F--> 3
4 --(--> 4
4 --id--> 5
6 --T--> 9
6 --F--> 3
6 --(--> 4
6 --id--> 5
7 --F--> 10
7 --(--> 4
7 --id--> 5
8 --+--> 6
8 --)--> 11
9 --*--> 7
========================================
FOLLOW Sets
========================================
FOLLOW(E) = { +, ), $ }
FOLLOW(T) = { +, *, ), $ }
FOLLOW(F) = { +, *, ), $ }
FOLLOW(E') = { $ }

========================================
ACTION Table
========================================
State   +       *       (       )       id      $       
0                       s4              s5              
1       s6                                      acc     
2       r2      s7              r2              r2      
3       r4      r4              r4              r4      
4                       s4              s5              
5       r6      r6              r6              r6      
6                       s4              s5              
7                       s4              s5              
8       s6                      s11                     
9       r1      s7              r1              r1      
10      r3      r3              r3              r3      
11      r5      r5              r5              r5      

========================================
GOTO Table
========================================
State   E       T       F       
0       1       2       3       
1                               
2                               
3                               
4       8       2       3       
5                               
6               9       3       
7                       10      
8                               
9                               
10                              
11                              

No SLR(1) conflicts detected.
\end{verbatim}
\textbf{说明}：该输出片段来自本次运行的完整输出文件（\texttt{output.txt}）。其中 FOLLOW 集用于归约项填表，ACTION 表显示各状态对终结符的移进/归约/接受动作。最后的 ``No SLR(1) conflicts detected'' 说明该文法在 SLR(1) 下无冲突。

\section{总结}
本实验成功完成了 SLR(1) 分析表的构造：
\begin{itemize}
	\item 基于 LR(0) 项目集规范族完成了状态转移图的构造。
	\item 计算得到 FIRST/FOLLOW 集并用于 SLR(1) 表填充。
	\item 最终生成 ACTION/GOTO 表，且对表达式文法未检测到冲突。
\end{itemize}
后续可扩展：增加对更多文法示例的测试，或进一步支持 LALR(1) 合并状态的分析表生成。

\section*{附录：实验源代码}
\subsection*{A.1 主要源码：main.c}
\begin{lstlisting}[language=C]
#include <ctype.h>
#include <stdio.h>
#include <stdlib.h>
#include <string.h>

#define MAX_LINE 1024
#define MAX_LINES 512
#define MAX_PRODS 512
#define MAX_RHS 32
#define MAX_SYMBOLS 256
#define MAX_STATES 256
#define MAX_ITEMS 4096
#define MAX_NAME 64

typedef struct
{
	int lhs;
	int rhs[MAX_RHS];
	int rhs_len;
} Production;

typedef struct
{
	int prod;
	int dot;
} Item;

typedef struct
{
	Item items[MAX_ITEMS];
	int count;
} ItemSet;

typedef enum
{
	ACT_NONE = 0,
	ACT_SHIFT,
	ACT_REDUCE,
	ACT_ACCEPT
} ActionType;

typedef struct
{
	ActionType type;
	int value;
} ActionCell;

static void
trim (char *s)
{
	size_t n = strlen (s);
	while (n > 0 && (s[n - 1] == '\n' || s[n - 1] == '\r'))
		{
			s[n - 1] = '\0';
			--n;
		}
	size_t i = 0;
	while (s[i] && isspace ((unsigned char)s[i]))
		{
			++i;
		}
	if (i > 0)
		{
			memmove (s, s + i, strlen (s + i) + 1);
		}
	n = strlen (s);
	while (n > 0 && isspace ((unsigned char)s[n - 1]))
		{
			s[n - 1] = '\0';
			--n;
		}
}

static int
is_blank (const char *s)
{
	while (*s)
		{
			if (!isspace ((unsigned char)*s))
				{
					return 0;
				}
			++s;
		}
	return 1;
}

static int
starts_with (const char *s, const char *prefix)
{
	return strncmp (s, prefix, strlen (prefix)) == 0;
}

static int
find_arrow (const char *s)
{
	const char *p = strstr (s, "->");
	if (!p)
		{
			return -1;
		}
	return (int)(p - s);
}

static int
name_index (char names[][MAX_NAME], int count, const char *name)
{
	for (int i = 0; i < count; ++i)
		{
			if (strcmp (names[i], name) == 0)
				{
					return i;
				}
		}
	return -1;
}

static int
add_name (char names[][MAX_NAME], int *count, const char *name)
{
	int idx = name_index (names, *count, name);
	if (idx >= 0)
		{
			return idx;
		}
	if (*count >= MAX_SYMBOLS)
		{
			return -1;
		}
	strncpy (names[*count], name, MAX_NAME - 1);
	names[*count][MAX_NAME - 1] = '\0';
	(*count)++;
	return (*count) - 1;
}

static int
is_nonterminal (char nonterms[][MAX_NAME], int nonterm_count, const char *s)
{
	return name_index (nonterms, nonterm_count, s) >= 0;
}

static int
item_equal (const Item *a, const Item *b)
{
	return a->prod == b->prod && a->dot == b->dot;
}

static int
itemset_contains (const ItemSet *set, const Item *it)
{
	for (int i = 0; i < set->count; ++i)
		{
			if (item_equal (&set->items[i], it))
				{
					return 1;
				}
		}
	return 0;
}

static void
itemset_add (ItemSet *set, const Item *it)
{
	if (!itemset_contains (set, it))
		{
			if (set->count < MAX_ITEMS)
				{
					set->items[set->count++] = *it;
				}
		}
}

static int
item_cmp (const void *a, const void *b)
{
	const Item *ia = (const Item *)a;
	const Item *ib = (const Item *)b;
	if (ia->prod != ib->prod)
		{
			return ia->prod - ib->prod;
		}
	return ia->dot - ib->dot;
}

static void
itemset_sort (ItemSet *set)
{
	qsort (set->items, (size_t)set->count, sizeof (Item), item_cmp);
}

static int
itemset_equal (const ItemSet *a, const ItemSet *b)
{
	if (a->count != b->count)
		{
			return 0;
		}
	for (int i = 0; i < a->count; ++i)
		{
			if (!item_equal (&a->items[i], &b->items[i]))
				{
					return 0;
				}
		}
	return 1;
}

static void
closure (ItemSet *set, const Production *prods, const int *nonterm_first,
				 const int *nonterm_count, const int *symbol_is_nonterm)
{
	int changed = 1;
	while (changed)
		{
			changed = 0;
			for (int i = 0; i < set->count; ++i)
				{
					Item it = set->items[i];
					const Production *p = &prods[it.prod];
					if (it.dot >= p->rhs_len)
						{
							continue;
						}
					int sym = p->rhs[it.dot];
					if (!symbol_is_nonterm[sym])
						{
							continue;
						}
					int nt_index = sym;
					int start = nonterm_first[nt_index];
					int cnt = nonterm_count[nt_index];
					if (start < 0 || cnt <= 0)
						{
							continue;
						}
					for (int k = 0; k < cnt; ++k)
						{
							Item new_item;
							new_item.prod = start + k;
							new_item.dot = 0;
							if (!itemset_contains (set, &new_item))
								{
									itemset_add (set, &new_item);
									changed = 1;
								}
						}
				}
		}
}

static void
goto_set (const ItemSet *src, int symbol, ItemSet *dst,
					const Production *prods, int prod_count, const int *nonterm_first,
					const int *nonterm_count, int nonterm_total,
					const int *symbol_is_nonterm)
{
	dst->count = 0;
	for (int i = 0; i < src->count; ++i)
		{
			Item it = src->items[i];
			const Production *p = &prods[it.prod];
			if (it.dot < p->rhs_len && p->rhs[it.dot] == symbol)
				{
					Item moved;
					moved.prod = it.prod;
					moved.dot = it.dot + 1;
					itemset_add (dst, &moved);
				}
		}
	(void)prod_count;
	(void)nonterm_total;
	closure (dst, prods, nonterm_first, nonterm_count, symbol_is_nonterm);
	itemset_sort (dst);
}

static void
print_production (const Production *p, char symbols[][MAX_NAME])
{
	printf ("%s ->", symbols[p->lhs]);
	if (p->rhs_len == 0)
		{
			printf (" epsilon");
		}
	else
		{
			for (int i = 0; i < p->rhs_len; ++i)
				{
					printf (" %s", symbols[p->rhs[i]]);
				}
		}
}

static void
print_item (const Item *it, const Production *prods, char symbols[][MAX_NAME])
{
	const Production *p = &prods[it->prod];
	printf ("%s ->", symbols[p->lhs]);
	for (int i = 0; i < p->rhs_len; ++i)
		{
			if (i == it->dot)
				{
					printf (" .");
				}
			printf (" %s", symbols[p->rhs[i]]);
		}
	if (it->dot == p->rhs_len)
		{
			printf (" .");
		}
}

static void
collect_kernel (const ItemSet *set, ItemSet *kernel)
{
	kernel->count = 0;
	for (int i = 0; i < set->count; ++i)
		{
			Item it = set->items[i];
			if (it.dot > 0 || it.prod == 0)
				{
					itemset_add (kernel, &it);
				}
		}
	itemset_sort (kernel);
}

static int
parse_grammar_lines (char lines[][MAX_LINE], int line_count,
										 char nonterms[][MAX_NAME], int *nonterm_count,
										 char symbols[][MAX_NAME], int *symbol_count,
										 int *start_symbol, Production *prods, int *prod_count)
{
	*nonterm_count = 0;
	*symbol_count = 0;
	*prod_count = 0;

	for (int i = 0; i < line_count; ++i)
		{
			char buf[MAX_LINE];
			strncpy (buf, lines[i], sizeof (buf) - 1);
			buf[sizeof (buf) - 1] = '\0';
			trim (buf);
			if (is_blank (buf) || starts_with (buf, "#") || starts_with (buf, "//"))
				{
					continue;
				}
			int arrow = find_arrow (buf);
			if (arrow < 0)
				{
					return 0;
				}
			buf[arrow] = '\0';
			buf[arrow + 1] = '\0';
			char lhs[MAX_NAME];
			strncpy (lhs, buf, sizeof (lhs) - 1);
			lhs[sizeof (lhs) - 1] = '\0';
			trim (lhs);
			if (lhs[0] == '\0')
				{
					return 0;
				}
			add_name (nonterms, nonterm_count, lhs);
		}

	if (*nonterm_count == 0)
		{
			return 0;
		}

	for (int i = 0; i < *nonterm_count; ++i)
		{
			add_name (symbols, symbol_count, nonterms[i]);
		}

	*start_symbol = 0;

	for (int i = 0; i < line_count; ++i)
		{
			char buf[MAX_LINE];
			strncpy (buf, lines[i], sizeof (buf) - 1);
			buf[sizeof (buf) - 1] = '\0';
			trim (buf);
			if (is_blank (buf) || starts_with (buf, "#") || starts_with (buf, "//"))
				{
					continue;
				}

			int arrow = find_arrow (buf);
			if (arrow < 0)
				{
					return 0;
				}

			buf[arrow] = '\0';
			char rhs_part[MAX_LINE];
			strncpy (rhs_part, buf + arrow + 2, sizeof (rhs_part) - 1);
			rhs_part[sizeof (rhs_part) - 1] = '\0';

			char lhs[MAX_NAME];
			strncpy (lhs, buf, sizeof (lhs) - 1);
			lhs[sizeof (lhs) - 1] = '\0';
			trim (lhs);
			if (lhs[0] == '\0')
				{
					return 0;
				}

			int lhs_sym = name_index (symbols, *symbol_count, lhs);
			if (lhs_sym < 0)
				{
					return 0;
				}

			char *alt = rhs_part;
			while (alt)
				{
					char *bar = strchr (alt, '|');
					if (bar)
						{
							*bar = '\0';
						}
					char alt_buf[MAX_LINE];
					strncpy (alt_buf, alt, sizeof (alt_buf) - 1);
					alt_buf[sizeof (alt_buf) - 1] = '\0';
					trim (alt_buf);

					Production p;
					p.lhs = lhs_sym;
					p.rhs_len = 0;

					if (alt_buf[0] == '\0' || strcmp (alt_buf, "epsilon") == 0
							|| strcmp (alt_buf, "eps") == 0)
						{
							p.rhs_len = 0;
						}
					else
						{
							char *tok = strtok (alt_buf, " \t");
							while (tok)
								{
									int sym;
									if (is_nonterminal (nonterms, *nonterm_count, tok))
										{
											sym = name_index (symbols, *symbol_count, tok);
										}
									else
										{
											sym = name_index (symbols, *symbol_count, tok);
											if (sym < 0)
												{
													sym = add_name (symbols, symbol_count, tok);
												}
										}
									if (sym < 0)
										{
											return 0;
										}
									if (p.rhs_len >= MAX_RHS)
										{
											return 0;
										}
									p.rhs[p.rhs_len++] = sym;
									tok = strtok (NULL, " \t");
								}
						}

					if (*prod_count >= MAX_PRODS)
						{
							return 0;
						}
					prods[(*prod_count)++] = p;

					if (bar)
						{
							alt = bar + 1;
						}
					else
						{
							alt = NULL;
						}
				}
		}

	return 1;
}

static int
add_to_set (unsigned char *set, int row, int col)
{
	unsigned char *cell = &set[row * MAX_SYMBOLS + col];
	if (*cell)
		{
			return 0;
		}
	*cell = 1;
	return 1;
}

static void
compute_first (const Production *prods, int prod_count, int symbol_count,
							 const int *symbol_is_nonterm, unsigned char *first,
							 unsigned char *first_eps)
{
	for (int i = 0; i < symbol_count; ++i)
		{
			first_eps[i] = 0;
			if (!symbol_is_nonterm[i])
				{
					add_to_set (first, i, i);
				}
		}

	int changed = 1;
	while (changed)
		{
			changed = 0;
			for (int i = 0; i < prod_count; ++i)
				{
					const Production *p = &prods[i];
					if (p->rhs_len == 0)
						{
							if (!first_eps[p->lhs])
								{
									first_eps[p->lhs] = 1;
									changed = 1;
								}
							continue;
						}

					int all_eps = 1;
					for (int j = 0; j < p->rhs_len; ++j)
						{
							int sym = p->rhs[j];
							for (int t = 0; t < symbol_count; ++t)
								{
									if (first[sym * MAX_SYMBOLS + t])
										{
											if (add_to_set (first, p->lhs, t))
												{
													changed = 1;
												}
										}
								}
							if (!first_eps[sym])
								{
									all_eps = 0;
									break;
								}
						}
					if (all_eps && !first_eps[p->lhs])
						{
							first_eps[p->lhs] = 1;
							changed = 1;
						}
				}
		}
}

static void
compute_follow (const Production *prods, int prod_count, int symbol_count,
								const int *symbol_is_nonterm, const unsigned char *first,
								const unsigned char *first_eps, unsigned char *follow,
								unsigned char *follow_end, int start_symbol)
{
	memset (follow_end, 0, MAX_SYMBOLS * sizeof (unsigned char));
	follow_end[start_symbol] = 1;

	int changed = 1;
	while (changed)
		{
			changed = 0;
			for (int i = 0; i < prod_count; ++i)
				{
					const Production *p = &prods[i];
					for (int j = 0; j < p->rhs_len; ++j)
						{
							int sym = p->rhs[j];
							if (!symbol_is_nonterm[sym])
								{
									continue;
								}

							int beta_has_eps = 1;
							for (int k = j + 1; k < p->rhs_len; ++k)
								{
									int next = p->rhs[k];
									for (int t = 0; t < symbol_count; ++t)
										{
											if (first[next * MAX_SYMBOLS + t])
												{
													if (add_to_set (follow, sym, t))
														{
															changed = 1;
														}
												}
										}
									if (!first_eps[next])
										{
											beta_has_eps = 0;
											break;
										}
								}

							if (j + 1 >= p->rhs_len)
								{
									beta_has_eps = 1;
								}

							if (beta_has_eps)
								{
									for (int t = 0; t < symbol_count; ++t)
										{
											if (follow[p->lhs * MAX_SYMBOLS + t])
												{
													if (add_to_set (follow, sym, t))
														{
															changed = 1;
														}
												}
										}
									if (follow_end[p->lhs] && !follow_end[sym])
										{
											follow_end[sym] = 1;
											changed = 1;
										}
								}
						}
				}
		}
}

static const char *
action_type_name (ActionType type)
{
	switch (type)
		{
		case ACT_SHIFT:
			return "shift";
		case ACT_REDUCE:
			return "reduce";
		case ACT_ACCEPT:
			return "accept";
		default:
			return "none";
		}
}

static int
set_action (ActionCell *cell, ActionType type, int value, int state,
						const char *sym)
{
	if (cell->type != ACT_NONE && (cell->type != type || cell->value != value))
		{
			printf ("Conflict at state %d, symbol %s: %s vs %s\n", state, sym,
							action_type_name (cell->type), action_type_name (type));
			return 1;
		}
	cell->type = type;
	cell->value = value;
	return 0;
}

static void
format_action (const ActionCell *cell, char *buf, size_t size)
{
	if (cell->type == ACT_SHIFT)
		{
			snprintf (buf, size, "s%d", cell->value);
		}
	else if (cell->type == ACT_REDUCE)
		{
			snprintf (buf, size, "r%d", cell->value);
		}
	else if (cell->type == ACT_ACCEPT)
		{
			snprintf (buf, size, "acc");
		}
	else
		{
			buf[0] = '\0';
		}
}

static void
print_follow_sets (int symbol_count, const int *symbol_is_nonterm,
									 char symbols[][MAX_NAME], const unsigned char *follow,
									 const unsigned char *follow_end)
{
	printf ("========================================\n");
	printf ("FOLLOW Sets\n");
	printf ("========================================\n");
	for (int i = 0; i < symbol_count; ++i)
		{
			if (!symbol_is_nonterm[i])
				{
					continue;
				}
			printf ("FOLLOW(%s) = { ", symbols[i]);
			int first_item = 1;
			for (int t = 0; t < symbol_count; ++t)
				{
					if (follow[i * MAX_SYMBOLS + t])
						{
							if (!first_item)
								{
									printf (", ");
								}
							printf ("%s", symbols[t]);
							first_item = 0;
						}
				}
			if (follow_end[i])
				{
					if (!first_item)
						{
							printf (", ");
						}
					printf ("$");
				}
			printf (" }\n");
		}
	printf ("\n");
}

static void
print_action_table (int state_count, const int *terminals, int term_count,
										char symbols[][MAX_NAME], const ActionCell *action_cells)
{
	printf ("========================================\n");
	printf ("ACTION Table\n");
	printf ("========================================\n");
	printf ("%-8s", "State");
	for (int i = 0; i < term_count; ++i)
		{
			printf ("%-8s", symbols[terminals[i]]);
		}
	printf ("%-8s\n", "$");

	for (int s = 0; s < state_count; ++s)
		{
			printf ("%-8d", s);
			for (int i = 0; i < term_count + 1; ++i)
				{
					const ActionCell *cell = &action_cells[s * (MAX_SYMBOLS + 1) + i];
					char buf[16];
					format_action (cell, buf, sizeof (buf));
					printf ("%-8s", buf);
				}
			printf ("\n");
		}
	printf ("\n");
}

static void
print_goto_table (int state_count, const int *nonterms, int nonterm_count,
									char symbols[][MAX_NAME], const int *goto_table)
{
	printf ("========================================\n");
	printf ("GOTO Table\n");
	printf ("========================================\n");
	printf ("%-8s", "State");
	for (int i = 0; i < nonterm_count; ++i)
		{
			printf ("%-8s", symbols[nonterms[i]]);
		}
	printf ("\n");

	for (int s = 0; s < state_count; ++s)
		{
			printf ("%-8d", s);
			for (int i = 0; i < nonterm_count; ++i)
				{
					int val = goto_table[s * MAX_SYMBOLS + i];
					if (val >= 0)
						{
							printf ("%-8d", val);
						}
					else
						{
							printf ("%-8s", "");
						}
				}
			printf ("\n");
		}
	printf ("\n");
}

int
main (int argc, char **argv)
{
	char (*lines)[MAX_LINE] = malloc (sizeof (*lines) * MAX_LINES);
	if (!lines)
		{
			fprintf (stderr, "Out of memory.\n");
			return 1;
		}
	int line_count = 0;

	FILE *fp = stdin;
	if (argc >= 2)
		{
			fp = fopen (argv[1], "r");
			if (!fp)
				{
					fprintf (stderr, "Cannot open file: %s\n", argv[1]);
					return 1;
				}
		}

	while (line_count < MAX_LINES && fgets (lines[line_count], MAX_LINE, fp))
		{
			trim (lines[line_count]);
			if (!is_blank (lines[line_count]))
				{
					line_count++;
				}
		}

	if (fp != stdin)
		{
			fclose (fp);
		}

	if (line_count == 0)
		{
			fprintf (stderr, "No grammar rules provided.\n");
			return 1;
		}

	char (*nonterms)[MAX_NAME] = malloc (sizeof (*nonterms) * MAX_SYMBOLS);
	char (*symbols)[MAX_NAME] = malloc (sizeof (*symbols) * MAX_SYMBOLS);
	if (!nonterms || !symbols)
		{
			fprintf (stderr, "Out of memory.\n");
			return 1;
		}
	int nonterm_count = 0;
	int symbol_count = 0;
	int start_symbol = 0;
	Production *prods = malloc (sizeof (*prods) * MAX_PRODS);
	if (!prods)
		{
			fprintf (stderr, "Out of memory.\n");
			return 1;
		}
	int prod_count = 0;

	if (!parse_grammar_lines (lines, line_count, nonterms, &nonterm_count,
														symbols, &symbol_count, &start_symbol, prods,
														&prod_count))
		{
			fprintf (stderr,
							 "Failed to parse grammar. Ensure lines like: E -> E + T | T\n");
			return 1;
		}

	if (argc >= 3)
		{
			int idx = name_index (symbols, symbol_count, argv[2]);
			if (idx < 0)
				{
					fprintf (stderr, "Start symbol '%s' not found in grammar.\n",
									 argv[2]);
					return 1;
				}
			start_symbol = idx;
		}

	char aug_name[MAX_NAME];
	snprintf (aug_name, sizeof (aug_name), "%s'", symbols[start_symbol]);
	if (name_index (symbols, symbol_count, aug_name) >= 0)
		{
			snprintf (aug_name, sizeof (aug_name), "%s0", symbols[start_symbol]);
		}

	int aug_sym = add_name (symbols, &symbol_count, aug_name);
	if (aug_sym < 0)
		{
			fprintf (stderr, "Too many symbols.\n");
			return 1;
		}

	Production aug;
	aug.lhs = aug_sym;
	aug.rhs_len = 1;
	aug.rhs[0] = start_symbol;

	Production all_prods[MAX_PRODS];
	int all_prod_count = 0;
	all_prods[all_prod_count++] = aug;
	for (int i = 0; i < prod_count; ++i)
		{
			all_prods[all_prod_count++] = prods[i];
		}

	int *symbol_is_nonterm = calloc (MAX_SYMBOLS, sizeof (int));
	if (!symbol_is_nonterm)
		{
			fprintf (stderr, "Out of memory.\n");
			return 1;
		}
	for (int i = 0; i < nonterm_count; ++i)
		{
			int idx = name_index (symbols, symbol_count, nonterms[i]);
			if (idx >= 0)
				{
					symbol_is_nonterm[idx] = 1;
				}
		}
	symbol_is_nonterm[aug_sym] = 1;

	int *nonterm_first = malloc (sizeof (int) * MAX_SYMBOLS);
	int *nonterm_prod_count = malloc (sizeof (int) * MAX_SYMBOLS);
	if (!nonterm_first || !nonterm_prod_count)
		{
			fprintf (stderr, "Out of memory.\n");
			return 1;
		}
	for (int i = 0; i < MAX_SYMBOLS; ++i)
		{
			nonterm_first[i] = -1;
			nonterm_prod_count[i] = 0;
		}

	for (int i = 0; i < all_prod_count; ++i)
		{
			int lhs = all_prods[i].lhs;
			if (nonterm_first[lhs] < 0)
				{
					nonterm_first[lhs] = i;
				}
			nonterm_prod_count[lhs]++;
		}

	for (int i = 0; i < symbol_count; ++i)
		{
			if (symbol_is_nonterm[i] && nonterm_prod_count[i] == 0)
				{
					fprintf (stderr, "Nonterminal '%s' has no productions.\n",
									 symbols[i]);
					return 1;
				}
		}

	printf ("========================================\n");
	printf ("Augmented Grammar\n");
	printf ("========================================\n");
	for (int i = 0; i < all_prod_count; ++i)
		{
			printf ("%d: ", i);
			print_production (&all_prods[i], symbols);
			printf ("\n");
		}

	ItemSet *states = calloc (MAX_STATES, sizeof (ItemSet));
	if (!states)
		{
			fprintf (stderr, "Out of memory.\n");
			return 1;
		}
	int state_count = 0;
	int (*transitions)[MAX_SYMBOLS]
			= malloc (sizeof (*transitions) * MAX_STATES);
	if (!transitions)
		{
			fprintf (stderr, "Out of memory.\n");
			return 1;
		}
	for (int i = 0; i < MAX_STATES; ++i)
		{
			for (int j = 0; j < MAX_SYMBOLS; ++j)
				{
					transitions[i][j] = -1;
				}
		}

	ItemSet start_set;
	start_set.count = 0;
	Item start_item;
	start_item.prod = 0;
	start_item.dot = 0;
	itemset_add (&start_set, &start_item);

	closure (&start_set, all_prods, nonterm_first, nonterm_prod_count,
					 symbol_is_nonterm);
	itemset_sort (&start_set);
	states[state_count++] = start_set;

	for (int i = 0; i < state_count; ++i)
		{
			for (int sym = 0; sym < symbol_count; ++sym)
				{
					ItemSet next;
					goto_set (&states[i], sym, &next, all_prods, all_prod_count,
										nonterm_first, nonterm_prod_count, symbol_count,
										symbol_is_nonterm);
					if (next.count == 0)
						{
							continue;
						}
					int existing = -1;
					for (int k = 0; k < state_count; ++k)
						{
							if (itemset_equal (&states[k], &next))
								{
									existing = k;
									break;
								}
						}
					if (existing < 0)
						{
							if (state_count >= MAX_STATES)
								{
									fprintf (stderr, "Too many states.\n");
									return 1;
								}
							states[state_count] = next;
							existing = state_count;
							state_count++;
						}
					transitions[i][sym] = existing;
				}
		}

	printf ("\n========================================\n");
	printf ("LR(0) Canonical Collection + Closure\n");
	printf ("========================================\n");

	for (int i = 0; i < state_count; ++i)
		{
			printf ("State %d:\n", i);
			for (int j = 0; j < states[i].count; ++j)
				{
					printf ("  ");
					print_item (&states[i].items[j], all_prods, symbols);
					printf ("\n");
				}
			ItemSet kernel;
			collect_kernel (&states[i], &kernel);
			printf ("Kernel: ");
			for (int k = 0; k < kernel.count; ++k)
				{
					if (k > 0)
						{
							printf (", ");
						}
					print_item (&kernel.items[k], all_prods, symbols);
				}
			printf ("\n\n");
		}

	printf ("========================================\n");
	printf ("State Transition Graph\n");
	printf ("========================================\n");
	for (int i = 0; i < state_count; ++i)
		{
			for (int sym = 0; sym < symbol_count; ++sym)
				{
					int to = transitions[i][sym];
					if (to >= 0)
						{
							printf ("%d --%s--> %d\n", i, symbols[sym], to);
						}
				}
		}

	unsigned char *first
			= calloc (MAX_SYMBOLS * MAX_SYMBOLS, sizeof (unsigned char));
	unsigned char *follow
			= calloc (MAX_SYMBOLS * MAX_SYMBOLS, sizeof (unsigned char));
	unsigned char *first_eps = calloc (MAX_SYMBOLS, sizeof (unsigned char));
	unsigned char *follow_end = calloc (MAX_SYMBOLS, sizeof (unsigned char));
	if (!first || !follow || !first_eps || !follow_end)
		{
			fprintf (stderr, "Out of memory.\n");
			return 1;
		}

	compute_first (all_prods, all_prod_count, symbol_count, symbol_is_nonterm,
								 first, first_eps);
	compute_follow (all_prods, all_prod_count, symbol_count, symbol_is_nonterm,
									first, first_eps, follow, follow_end, aug_sym);
	print_follow_sets (symbol_count, symbol_is_nonterm, symbols, follow,
										 follow_end);

	int terminals[MAX_SYMBOLS];
	int term_count = 0;
	int term_col[MAX_SYMBOLS];
	for (int i = 0; i < MAX_SYMBOLS; ++i)
		{
			term_col[i] = -1;
		}
	for (int i = 0; i < symbol_count; ++i)
		{
			if (!symbol_is_nonterm[i])
				{
					term_col[i] = term_count;
					terminals[term_count++] = i;
				}
		}

	int nonterms_out[MAX_SYMBOLS];
	int nonterm_out_count = 0;
	for (int i = 0; i < symbol_count; ++i)
		{
			if (symbol_is_nonterm[i] && i != aug_sym)
				{
					nonterms_out[nonterm_out_count++] = i;
				}
		}

	ActionCell *action
			= calloc (MAX_STATES * (MAX_SYMBOLS + 1), sizeof (ActionCell));
	int *goto_table = malloc (sizeof (int) * MAX_STATES * MAX_SYMBOLS);
	if (!action || !goto_table)
		{
			fprintf (stderr, "Out of memory.\n");
			return 1;
		}
	for (int i = 0; i < MAX_STATES * MAX_SYMBOLS; ++i)
		{
			goto_table[i] = -1;
		}

	int conflict = 0;
	for (int i = 0; i < state_count; ++i)
		{
			for (int j = 0; j < states[i].count; ++j)
				{
					Item it = states[i].items[j];
					const Production *p = &all_prods[it.prod];
					if (it.dot < p->rhs_len)
						{
							int sym = p->rhs[it.dot];
							if (!symbol_is_nonterm[sym])
								{
									int to = transitions[i][sym];
									int col = term_col[sym];
									if (to >= 0 && col >= 0)
										{
											ActionCell *cell = &action[i * (MAX_SYMBOLS + 1) + col];
											if (set_action (cell, ACT_SHIFT, to, i, symbols[sym]))
												{
													conflict = 1;
												}
										}
								}
						}
					else
						{
							if (it.prod == 0)
								{
									ActionCell *cell
											= &action[i * (MAX_SYMBOLS + 1) + term_count];
									if (set_action (cell, ACT_ACCEPT, 0, i, "$"))
										{
											conflict = 1;
										}
								}
							else
								{
									int lhs = p->lhs;
									for (int t = 0; t < symbol_count; ++t)
										{
											if (follow[lhs * MAX_SYMBOLS + t])
												{
													int col = term_col[t];
													if (col >= 0)
														{
															ActionCell *cell
																	= &action[i * (MAX_SYMBOLS + 1) + col];
															if (set_action (cell, ACT_REDUCE, it.prod, i,
																							symbols[t]))
																{
																	conflict = 1;
																}
														}
												}
										}
									if (follow_end[lhs])
										{
											ActionCell *cell
													= &action[i * (MAX_SYMBOLS + 1) + term_count];
											if (set_action (cell, ACT_REDUCE, it.prod, i, "$"))
												{
													conflict = 1;
												}
										}
								}
						}
				}

			for (int j = 0; j < nonterm_out_count; ++j)
				{
					int sym = nonterms_out[j];
					int to = transitions[i][sym];
					if (to >= 0)
						{
							goto_table[i * MAX_SYMBOLS + j] = to;
						}
				}
		}

	print_action_table (state_count, terminals, term_count, symbols, action);
	print_goto_table (state_count, nonterms_out, nonterm_out_count, symbols,
										goto_table);

	if (conflict)
		{
			printf ("Grammar has SLR(1) conflicts.\n");
		}
	else
		{
			printf ("No SLR(1) conflicts detected.\n");
		}

	return 0;
}
\end{lstlisting}

\subsection*{A.2 运行命令}
\begin{verbatim}
# 编译
gcc -Wall -Wextra -o slr1 main.c

# 运行（表达式文法）
./slr1 grammar_expr.txt E
\end{verbatim}

\end{document}
